% ============================================================
%  CHAPTER 11 — Discussion
% ============================================================
\chapter{Discussion}
\label{ch:discussion}

\section{Key Findings}

\subsection{VLM-Based Task Decomposition}

The use of Gemini 2.0 Flash as the reasoning engine proved effective for multi-step
task decomposition.
In successful trials, the VLM consistently selected the appropriate skill sequence
without explicit programming of the skill ordering.
The structured JSON output format, with the system prompt constraint
\textit{``Respond ONLY with valid JSON''}, reduced parse failures to below \todo{X}\%
of queries, with the retry logic handling the remainder.

The affordance feedback loop — where skill failures are communicated back to the
VLM as context — was essential for robustness.
Without it, the VLM would repeatedly select an infeasible skill; with it, the agent
dynamically adapts, for example switching from \texttt{grab\_bottle} to
\texttt{search\_with\_head} when the object is not visible.

\subsection{3-D Localisation Accuracy}

The RANSAC cylinder fitting approach provided stable object localisation with typical
standard deviations below 3 mm across 5 frames.
This accuracy was sufficient for the gripper to close around the bottle reliably,
given the calibrated grasp offsets.
The main failure mode was insufficient depth data — thin objects (\eg a narrow bottle
at the edge of the YOLO bounding box) produced fewer than 15 valid depth points,
causing the cylinder fit to be rejected.

\subsection{Torso Descent vs.\ Cartesian Approach}

The torso descent strategy was an unconventional but effective solution to the wrist
singularity problem that caused MoveIt Cartesian path failures at low gripper heights.
By using the prismatic torso joint (world-Z aligned) for the final vertical approach,
the system guaranteed a straight-line end-effector path without requiring IK solutions
at each intermediate height.

The main drawback is speed: the torso descends at 3 cm/s for safety, making the
approach phase 3--4 seconds regardless of the required travel distance.
The planned MoveJ/MoveL trajectory approach (Chapter~\ref{ch:manipulation}) would
reduce this to 1--2 seconds while providing a more professional motion profile.

\subsection{Split-Process Architecture}

The HTTP microservice architecture for YOLO inference demonstrated that it is possible
to overcome severe deployment constraints (Python 3.6, no GPU) while maintaining
real-time performance.
The 18 ms YOLO11n inference time on the host CPU, combined with HTTP overhead, resulted
in an effective detection rate of \todo{X} Hz, which was adequate for the agent's
non-real-time control loop.

The main risk of this architecture is the dependency on a stable network connection
between the container and the host, and the need to start two separate processes.
The fallback to local YOLOv4-tiny mitigates service unavailability.

\section{Limitations}

\subsection{No Object Permanence}

Objects are tracked only within the current camera frame.
If the bottle is moved out of view during a grasp attempt (for example, by being
knocked), the system has no memory of its last known location.
The next perception step would find no bottle, and the VLM would select
\texttt{search\_with\_head}, which may or may not relocate it depending on where
it fell.

\subsection{Fixed Grasp Orientation}

The grasp orientation quaternion is fixed throughout the pipeline.
This works well for upright cylindrical objects on flat surfaces but would fail for:
\begin{itemize}[noitemsep]
  \item Objects lying on their side.
  \item Objects oriented at an angle to the robot's forward direction.
  \item Non-cylindrical objects requiring a different approach angle.
\end{itemize}

Adapting the grasp orientation to the detected object pose (using either
learning-based pose estimation or RANSAC-based axis detection) is an important
future direction.

\subsection{Internet Dependency}

The VLM reasoning requires internet access for the Gemini API.
Network interruptions degrade reasoning quality (fallback to \texttt{go\_home}) and
latency spikes (up to 1.5 s) slow the control loop.
A locally deployed VLM (LLaVA, VILA) would eliminate this dependency at the cost
of requiring a GPU.

\subsection{Python 3.6 Constraint}

The Docker container's Python 3.6 requirement prevented the use of several beneficial
libraries:
\begin{itemize}[noitemsep]
  \item \texttt{moveit\_commander} (Python 2.7 or 3.8+ only in the PAL stack).
  \item Modern type annotations (\texttt{x: int | None}).
  \item \texttt{dataclasses} (Python 3.7+).
\end{itemize}

The workarounds (raw MoveGroupAction client, manual type checking, dictionary-based
state) add code complexity.

\subsection{Single-Object Grasping}

The system can only grasp one object per task cycle.
Multi-object grasping (carrying a tray, for example) would require significant
extensions to the perception, planning, and skill layers.

\section{Comparison with Related Work}

\begin{table}[H]
\centering
\caption{Comparison with related VLA robot systems.}
\label{tab:related_comparison}
\begin{tabularx}{\linewidth}{lXXX}
\toprule
\textbf{System} & \textbf{VLM/LLM} & \textbf{Skills} & \textbf{Grasp approach} \\
\midrule
SayCan \cite{ahn2022saycan} & PaLM + value fn & 551 skills & Learned policy \\
RT-2 \cite{brohan2023rt2}   & PaLM-E (fine-tuned) & End-to-end & Learned trajectory \\
This work                    & Gemini 2.0 Flash (zero-shot) & 9 skills & Classical RANSAC + MoveIt \\
\bottomrule
\end{tabularx}
\end{table}

Unlike SayCan and RT-2, this work requires no robot-specific training data and can be
deployed on a new platform with only calibration of the grasp offsets.
The trade-off is reduced generalisation: the skills are hand-crafted for specific
manipulation primitives rather than learned from demonstrations.

\section{Lessons Learned}

\begin{enumerate}
  \item \textbf{Iterative development is essential.} The grasp pipeline required 5
        iterations to achieve reliable performance. Each version addressed a specific
        failure mode identified through real-robot testing.

  \item \textbf{Sensor noise must be actively managed.} Single-frame detection is
        too noisy for reliable 3-D localisation; 5-frame averaging reduced grasp
        failures significantly.

  \item \textbf{Structured prompts are critical.} Early prompt designs without strict
        JSON formatting constraints produced unparseable outputs frequently. The current
        design reduces failures to below \todo{X}\%.

  \item \textbf{Safety should be designed in from the start.} The affordance gating
        system prevented several potentially damaging arm movements during development.

  \item \textbf{Architecture constraints drive design.} The Python 3.6 constraint
        led to the HTTP microservice design, which turned out to be a clean separation
        of concerns that would be useful even without the constraint.
\end{enumerate}
